\chapter{P\&L explain et vérification indépendante des prix}\label{ch:rc:pnl-explain-and-independent-price-verification}

Chaque soir, un product controller doit expliquer chaque dollar qu'un desk a gagné ou perdu~: une part
vient du mouvement des taux, une autre de la volatilité, une autre du passage du temps, et il reste un
résidu que personne ne sait justifier. La plupart du temps, ce résidu est faible. Le jour où un portefeuille de swaptions a gagné
4.57 millions de dollars après une hausse des taux de 30 points de base et une remontée de la volatilité~;
l'explication propre au desk en couvrait 4.17 millions de dollars, et 400 milliers de dollars restaient inexpliqués. Un
\omterm{def:rc:pnl-explain-and-independent-price-verification:unexplained}{P\&L inexpliqué} traduit soit un risque manquant, soit une valorisation erronée, soit une erreur de comptabilisation, et le \omterm{def:rc:pnl-explain-and-independent-price-verification:pc}{product control}
doit déterminer laquelle. Ce chapitre présente les outils~: l'\omterm{def:rc:pnl-explain-and-independent-price-verification:risk}{explication par les sensibilités} et par réévaluation,
la vérification des valorisations par rapport à des prix indépendants, les réserves et les
corrections de valeur prudentielles qui tiennent compte de ce que les valorisations ne peuvent pas savoir, et les fonctions de contrôle
qui les mettent en œuvre. Le livre~1 a introduit l'attribution du P\&L pour une société de trading~;
ici, elle est appliquée à un portefeuille de produits dérivés, où les risques interagissent.

\section{Attribution~: par les sensibilités et par réévaluation}

\begin{definition}[Explication par les sensibilités]\label{def:rc:pnl-explain-and-independent-price-verification:risk}
L'\emph{explication par les sensibilités}\index{explication par les sensibilités} décompose le P\&L de la
journée par un développement de Taylor selon les variations des \omterm{def:rc:market-risk-measures:factor}{facteurs de risque}, en utilisant les grecques d'ouverture~:
$\Delta V\approx\delta\,\Delta f+\tfrac12\Gamma\,\Delta f^2+\nu\,\Delta\sigma+\tfrac12\,\mathrm{volga}\,\Delta\sigma^2+\mathrm{vanna}\,\Delta f\,\Delta\sigma+\theta\,\Delta t$.
\end{definition}

\begin{definition}[Explication par réévaluation]\label{def:rc:pnl-explain-and-independent-price-verification:reval}
L'\emph{explication par réévaluation}\index{explication par réévaluation} fait passer les \omterm{def:rc:market-risk-measures:factor}{facteurs de risque}
de leurs valeurs d'ouverture à leurs valeurs de clôture un à un, en réévaluant intégralement l'instrument
après chaque étape et affecte la variation de chaque étape à son facteur. Les étapes s'additionnent exactement au P\&L~;
leur répartition dépend toutefois de l'ordre choisi.
\end{definition}

\begin{definition}[P\&L inexpliqué]\label{def:rc:pnl-explain-and-independent-price-verification:unexplained}
Le \emph{P\&L inexpliqué}\index{P\&L inexpliqué} est l'écart entre le P\&L réel et la somme des
composantes attribuées. Les contrôleurs fixent des seuils sur ce montant, en valeur et en proportion du
P\&L, au-delà desquels il faut procéder à une investigation.
\end{definition}

\begin{example}[Un risk reversal lors d'une journée calme]\label{ex:rc:pnl-explain-and-independent-price-verification:book}
Un portefeuille illustratif est long de 320 millions de dollars de swaptions payeuses 1 an sur 10 ans, de prix d'exercice 4.50~\%,
et court du même notionnel de swaptions receveuses à 3.50~\%, soit un risk reversal autour d'un taux forward de 4.00~\%, valorisé à
une volatilité normale de 95 points de base (0.75 an avant l'échéance, annuité 8.0). Les deux jambes valent la
même chose, si bien que le portefeuille vaut zéro~; par symétrie, son gamma, son vega et son theta sont également nuls~: c'est une
position de pur delta, de 139\,103~dollars par point de base. Sa sensibilité croisée n'est pas nulle~: le vanna
vaut 1\,086~dollars par point de base de taux et par point de base de volatilité, car une hausse du taux rapproche la payeuse longue
de la monnaie et éloigne la receveuse courte de celle-ci.
\end{example}

\omcode{code/firm/pnlexplain/firm_pnlexplain.py}{19}{41}{Grecques par différences finies, y compris le terme croisé, et explication par les sensibilités avec ou sans ce terme.}\label{lst:rc:pnl-explain-and-independent-price-verification:risk}

\section{Explication du P\&L inexpliqué et ce qu'il révèle}

\begin{example}[La journée des grands mouvements]\label{ex:rc:pnl-explain-and-independent-price-verification:day}
Le forward monte de 30 points de base et la volatilité de 12. Le portefeuille gagne 4\,571\,501~dollars. L'explication
du desk, avec delta, gamma, vega et theta, donne 4\,173\,095~dollars, tout provenant du delta~: 398\,406~dollars restent inexpliqués,
soit 9~\% de la journée. En ajoutant le terme de vanna, $1\,086\times30\times12$, on explique
391\,087~dollars de ce montant~; il en reste 7\,319, dus aux termes d'ordre supérieur (la \cref{fig:rc:pnl-explain-and-independent-price-verification:attribution}). Le
\omterm{def:rc:pnl-explain-and-independent-price-verification:unexplained}{P\&L inexpliqué} croît avec le produit des mouvements de taux et de volatilité, ce qui explique qu'il apparaisse les grands jours
(la \cref{fig:rc:pnl-explain-and-independent-price-verification:unexplained}).
\end{example}

\begin{omfigure}
\begin{tikzpicture}
\begin{axis}[width=11.5cm,height=5.2cm,ybar=4pt,bar width=24pt,
  symbolic x coords={delta,vanna,unexplained},xtick=data,xticklabels={delta,vanna,inexpliqué},enlarge x limits=0.25,
  ylabel={milliers de USD},
  tick label style={font=\small},label style={font=\small},
  ymin=0,ymax=4800,grid=major,grid style={gray!25},
  nodes near coords,every node near coord/.append style={font=\footnotesize,/pgf/number format/fixed,/pgf/number format/precision=0},
  legend columns=2,legend style={font=\footnotesize,at={(0.5,-0.2)},anchor=north,draw=none,fill=none,/tikz/every even column/.append style={column sep=8pt}},
  table/col sep=comma]
\addplot[fill=omExo!40,draw=omExo] table[x=term,y=desk]{figdata/rates-credit-risk/27-pnl-explain-and-independent-price-verification/attribution.csv};
\addplot[fill=omDef!60,draw=omDef] table[x=term,y=full]{figdata/rates-credit-risk/27-pnl-explain-and-independent-price-verification/attribution.csv};
\legend{explication du desk,avec le terme croisé}
\end{axis}
\end{tikzpicture}

\omcaption{Attribution du P\&L du risk reversal lors de la journée des grands mouvements (4.57~millions de dollars)~: l'explication du desk
laisse 398~milliers de dollars inexpliqués~; le terme de vanna en couvre presque la totalité. Données~: le tutoriel
du chapitre.}\label{fig:rc:pnl-explain-and-independent-price-verification:attribution}
\end{omfigure}

\begin{omfigure}
\begin{tikzpicture}
\begin{axis}[width=11.5cm,height=5.2cm,
  xlabel={variation du taux forward (pb), volatilité en hausse de 12 pb},ylabel={inexpliqué (milliers de USD)},
  tick label style={font=\small},label style={font=\small},
  xmin=-40,xmax=40,ymin=-650,ymax=650,grid=major,grid style={gray!25},
  legend columns=2,legend style={font=\footnotesize,at={(0.5,-0.25)},anchor=north,draw=none,fill=none,/tikz/every even column/.append style={column sep=8pt}},
  table/col sep=comma]
\addplot[omThm,very thick,mark=*,mark size=1.2pt] table[x=bp,y=desk]{figdata/rates-credit-risk/27-pnl-explain-and-independent-price-verification/unexplained.csv};
\addplot[omDef,very thick,mark=square*,mark size=1.2pt] table[x=bp,y=full]{figdata/rates-credit-risk/27-pnl-explain-and-independent-price-verification/unexplained.csv};
\legend{explication du desk,avec le terme croisé}
\end{axis}
\end{tikzpicture}

\omcaption{\omterm{def:rc:pnl-explain-and-independent-price-verification:unexplained}{P\&L inexpliqué} du risk reversal en fonction de la variation de taux du jour, la volatilité étant en hausse
de 12 points de base. Sans le terme croisé, la part inexpliquée est proportionnelle à la variation de taux~; avec lui,
il ne subsiste que des termes d'ordre trois. Données~: le tutoriel du chapitre.}\label{fig:rc:pnl-explain-and-independent-price-verification:unexplained}
\end{omfigure}

\begin{example}[La même journée par réévaluation]\label{ex:rc:pnl-explain-and-independent-price-verification:reval}
En faisant passer d'abord le forward, puis la volatilité, puis le temps~: le forward contribue pour 4\,240\,447~dollars,
la volatilité pour 336\,046~dollars, le temps pour $-4\,992$. En faisant passer d'abord la volatilité~: la volatilité ne contribue pas (le
vega du portefeuille est nul au départ), le forward pour 4\,576\,493~dollars. Les deux résultats s'additionnent au P\&L~; l'effet croisé
revient au facteur qui bouge en second.
\end{example}

\begin{definition}[Product control]\label{def:rc:pnl-explain-and-independent-price-verification:pc}
Le \emph{product control}\index{product control} est la fonction financière, indépendante des desks, qui
produit et explique le P\&L quotidien, vérifie les valorisations, fixe les réserves et réconcilie les
\omterm{def:rc:regulatory-capital-for-trading-books:books}{portefeuilles de négociation} avec le grand livre.
\end{definition}

\section{Vérification indépendante des prix}

\begin{definition}[Vérification indépendante des prix]\label{def:rc:pnl-explain-and-independent-price-verification:ipv}
La \emph{vérification indépendante des prix}\index{vérification indépendante des prix} (IPV) est le processus par
lequel les prix de marché et les données d'entrée des modèles utilisés par un desk sont vérifiés par rapport à des sources indépendantes de lui
(cotations de courtiers, transactions, services de consensus) et les valorisations sont corrigées lorsqu'elles sortent d'une
tolérance.
\end{definition}

Dans un service de consensus, les dealers soumettent leurs prix milieu pour une liste de contrats une fois par mois et reçoivent
la moyenne des soumissions acceptées~; un dealer dont la soumission est rejetée ne reçoit pas le
consensus. Le service de référence, Totem, a débuté en 1997 avec six dealers et comptait environ 120 participants et quelque 30
contributeurs par contrat, selon une étude de 2020 portant sur ses données.

\omcode{code/firm/pnlexplain/firm_pnlexplain.py}{56}{62}{IPV~: une valorisation hors de la bande de consensus est ramenée à la borne la plus proche.}\label{lst:rc:pnl-explain-and-independent-price-verification:ipv}

\begin{example}[Vérification des valorisations de volatilité]\label{ex:rc:pnl-explain-and-independent-price-verification:ipv}
Le desk valorise les deux strikes à 95 points de base~; un consensus illustratif donne 97 pour la payeuse à 4.50~\%
et 91 pour la receveuse à 3.50~\%, soit le skew du marché, avec une tolérance de 2 points de base. La valorisation de la payeuse
se situe sur la borne et est conservée~; celle de la receveuse est ramenée à 93, ce qui diminue la valeur des receveuses vendues et ajoute 146\,482~dollars.
Les valorisations plates expliquent aussi la vanna~: un desk qui ne valorise aucun skew ne voit
aucune raison pour que son vega change lorsque les taux bougent.
\end{example}

\begin{dated}{2026-09}{Valorisations à la borne}\label{dat:rc:pnl-explain-and-independent-price-verification:cio}
Le rapport 2013 du groupe de travail de JPMorgan décrit les tests de prix de son Chief Investment Office~: chaque valorisation de trader
était comparée à un prix milieu calculé par le groupe de contrôle des valorisations et à un seuil reflétant
la fourchette cours acheteur--cours vendeur, et les valorisations situées hors de celle-ci devaient être ramenées à la borne la plus proche. Lors des tests de mars 2012,
de nombreuses positions étaient valorisées sur la borne de la fourchette cours acheteur--cours vendeur ou près d'elle~; se situant dans les
seuils, elles ont été acceptées.
\end{dated}

\section{Réserves et hiérarchie de la juste valeur}

\begin{definition}[Hiérarchie de la juste valeur]\label{def:rc:pnl-explain-and-independent-price-verification:hierarchy}
La \emph{hiérarchie de la juste valeur}\index{hiérarchie de la juste valeur} classe les positions évaluées à la juste valeur selon le
caractère observable de leurs données d'entrée (IFRS~13)~: niveau~1, cours cotés non ajustés sur des marchés actifs pour
des instruments identiques~; niveau~2, autres données d'entrée directement ou indirectement observables~; niveau~3, données
d'entrée non observables. Les positions de niveau~3 demandent le
plus de vérifications et portent les réserves les plus importantes.
\end{definition}

Les réserves corrigent les justes valeurs de ce qu'ignore un prix milieu~: le coût cours acheteur--cours vendeur de sortie,
l'incertitude du modèle et de ses paramètres, et les profits du premier jour qui ne peuvent pas encore être reconnus
(livre~5, chapitre~27). Elles font partie de la juste valeur~; les régulateurs en demandent davantage.

\section{Évaluation prudente}

\begin{definition}[Évaluation prudente, correction de valeur supplémentaire]\label{def:rc:pnl-explain-and-independent-price-verification:pv}
L'\emph{évaluation prudente}\index{evaluation@évaluation prudente} valorise les positions évaluées à la juste valeur au prix auquel
la banque a 90~\% de confiance de pouvoir les sortir. Les \emph{corrections de valeur
supplémentaires}\index{correction de valeur supplémentaire} (AVA) sont les écarts entre la juste valeur et
la valeur prudente, par catégorie (incertitude sur le prix de marché, coûts de clôture, \omterm{def:rc:model-risk-and-validation:risk}{risque de modèle} et autres)~; elles sont
déduites des fonds propres durs.
\end{definition}

\begin{example}[L'AVA du portefeuille]\label{ex:rc:pnl-explain-and-independent-price-verification:ava}
Avec des volatilités des contributeurs dispersées de 2.0 et 2.5 points de base autour du consensus, l'AVA d'incertitude
sur le prix de marché, soit la juste valeur moins la valeur de sortie à 90~\% de confiance, vaut 190\,578~dollars pour les payeuses et
232\,964 pour les receveuses~; sommées après le facteur d'agrégation de 50~\%, elles donnent 211\,771. L'approche simplifiée,
0.1~\% des valeurs absolues de juste valeur (6.87~millions de dollars) donnerait 6\,872~: pour un portefeuille d'options dont la valeur est
faible au regard de son risque, le chiffre simplifié sous-estime l'incertitude.
\end{example}

\begin{dated}{2026-09}{Évaluation prudente dans l'Union européenne}\label{dat:rc:pnl-explain-and-independent-price-verification:pv}
Selon la norme européenne d'\omterm{def:rc:pnl-explain-and-independent-price-verification:pv}{évaluation prudente} (Delegated Regulation 2016/101), la valeur prudente est la valeur
à laquelle une banque a 90~\% de confiance de pouvoir sortir~; les établissements dont la valeur absolue des actifs et passifs évalués à la juste valeur est inférieure à 15~milliards d'euros
peuvent utiliser l'approche simplifiée (0.1~\% de cette somme), les autres l'approche de base, dans laquelle
les AVA d'incertitude sur le prix de marché, de coût de clôture et de \omterm{def:rc:model-risk-and-validation:risk}{risque de modèle} sont sommées après un facteur d'agrégation
de 50~\% (temporairement porté à 66~\% après l'amendement de 2020 pour volatilité extrême). Les AVA totaux sont déduits des
fonds propres durs de catégorie~1.
\end{dated}

\section{Tutoriel~: expliquer une journée}\label{tut:rc:pnl-explain-and-independent-price-verification:tutorial}

\begin{tutorial}
\textbf{Objectif.} Expliquer la journée à grand mouvement du risk reversal par les grecques et par réévaluation, vérifier ses valorisations et
calculer son ajustement d'\omterm{def:rc:pnl-explain-and-independent-price-verification:pv}{évaluation prudente}. \textbf{État final~:} les chiffres des
\cref{ex:rc:pnl-explain-and-independent-price-verification:book,ex:rc:pnl-explain-and-independent-price-verification:day,ex:rc:pnl-explain-and-independent-price-verification:reval,ex:rc:pnl-explain-and-independent-price-verification:ipv,ex:rc:pnl-explain-and-independent-price-verification:ava}
et les graphiques.
\begin{tutsteps}
\item \textbf{Grecques}~: \texttt{greeks(value, START, BUMPS)}.
\item \textbf{Explication}~: \texttt{day()} avec et sans le terme croisé~; \texttt{revaluation()} selon deux ordres.
\item \textbf{IPV et AVA}~: \texttt{ipv\_table()}, \texttt{ava\_table()}.
\item \textbf{Graphiques}~: \texttt{fig\_rc\_pnlexplain.py}.
\end{tutsteps}
\textbf{Pour aller plus loin.} Valoriser le skew à partir du consensus et refaire la journée~: la vanna devient un vega
que l'explication perçoit~; ajouter un second forward et un terme croisé de corrélation.
\end{tutorial}

\section{Construction~: l'explication du P\&L}\label{bld:rc:pnl-explain-and-independent-price-verification:pnlexplain}

\begin{build}
\bfield{Rôle} L'explication quotidienne de la société, la vérification des prix et l'\omterm{def:rc:pnl-explain-and-independent-price-verification:pv}{évaluation prudente}, alimentées par le P\&L
du gestionnaire de positions du livre~1 (\texttt{firm.pnl}) et par les grecques de la bibliothèque de pricing.
\bfield{Interface} \texttt{greeks(value, state, bumps)}~; \texttt{risk\_based(g, df, ds, days, cross)}~;
\texttt{revaluation\_based(value, start, end, order)}~; \texttt{ipv(marks, consensus, tolerance)}~;
\texttt{mpu\_ava}~; \texttt{aggregate\_ava}~; \texttt{simplified\_ava}.
\bfield{Règles} Grecques en début de journée~; seuils d'explication exprimés en montant et en actions~; l'IPV ramène les valorisations à la
borne de tolérance~; confiance à 90~\% et agrégation à 50~\% pour les AVA.
\bfield{Tests d'acceptation} \texttt{code/firm/pnlexplain/tests/}~: un portefeuille quadratique est expliqué exactement avec le
terme croisé et pas sans lui~; les étapes de réévaluation s'additionnent et dépendent de l'ordre~; bornes de l'IPV~; AVA d'une
dispersion normale~; agrégation et chiffres simplifiés.
\bfield{Pour aller plus loin} Explication par transaction et par \omterm{def:rc:market-risk-measures:factor}{facteur de risque} sur un portefeuille entier~; alertes sur le \omterm{def:rc:pnl-explain-and-independent-price-verification:unexplained}{P\&L inexpliqué}~; coût de
clôture et AVA de \omterm{def:rc:model-risk-and-validation:risk}{risque de modèle}~; ingestion des données de consensus.
\end{build}

\begin{omsources}
\item European Banking Authority, final draft RTS amending the RTS on prudent valuation, EBA/RTS/2020/04, 2020.
\item Basel Committee on Banking Supervision, \emph{Supervisory guidance for assessing banks' financial
instrument fair value practices}, avril 2009.
\item JPMorgan Chase \& Co., \emph{Report of the Management Task Force Regarding 2012 CIO Losses}, 2013.
\item IFRS~13 \emph{Fair Value Measurement}, paragraphes 76, 81 et 86, tel qu'adopté dans
la Commission Regulation (EU) n°~1255/2012.
\item L.~M.\ Ergun et A.~Uthemann, «~Higher-order uncertainty in financial markets: evidence from a
consensus pricing service~», Systemic Risk Centre Discussion Paper 98, juin 2020.
\end{omsources}

\section{Exercices}

\begin{exercise}[$\star$]\label{exo:rc:pnl-explain-and-independent-price-verification:1}
Calculer le terme de vanna du jour de grand mouvement à partir de la vanna du portefeuille, égale à 1\,086~dollars par point de base au carré.
\end{exercise}

\begin{exercise}[$\star$]\label{exo:rc:pnl-explain-and-independent-price-verification:2}
Pourquoi les grecques gamma, vega et theta du risk reversal sont-elles nulles en début de journée~?
\end{exercise}

\begin{exercise}[$\star$]\label{exo:rc:pnl-explain-and-independent-price-verification:3}
Une valorisation est fixée à 38 contre un prix milieu consensuel de 35, avec une tolérance de 2. Quelle est la valorisation vérifiée~?
\end{exercise}

\begin{exercise}[$\star\star$]\label{exo:rc:pnl-explain-and-independent-price-verification:4}
Pourquoi l'\omterm{def:rc:pnl-explain-and-independent-price-verification:reval}{explication par réévaluation} dépend-elle de l'ordre des facteurs, et comment
supprimer cette dépendance~?
\end{exercise}

\begin{exercise}[$\star\star$]\label{exo:rc:pnl-explain-and-independent-price-verification:5}
Ramener le \omterm{def:rc:pnl-explain-and-independent-price-verification:unexplained}{P\&L inexpliqué} à une journée avec un mouvement de taux de 15~points de base et le même mouvement de volatilité.
\end{exercise}

\begin{exercise}[$\star\star$]\label{exo:rc:pnl-explain-and-independent-price-verification:6}
Pourquoi l'AVA d'incertitude sur le prix de marché d'un portefeuille d'options est-elle grande par rapport à sa fair value~?
\end{exercise}

\begin{exercise}[$\star\star\star$]\label{exo:rc:pnl-explain-and-independent-price-verification:7}
\emph{Programmation.} Calculer l'AVA du portefeuille lorsque la dispersion des contributeurs est doublée.
\end{exercise}

\begin{exercise}[$\star\star\star$]\label{exo:rc:pnl-explain-and-independent-price-verification:8}
\emph{Trouver la faille.} «~Toutes nos marques sont dans les seuils de l'IPV, notre valorisation est donc satisfaisante.~»
\end{exercise}

\section{Problème~: Les 400~000 inexpliqués}

\begin{problem}[{Problème du week-end --- le terme croisé manquant}]\label{pb:rc:pnl-explain-and-independent-price-verification:1}
Le product controller du desk de risk reversal constate 398\,406~USD inexpliqués le jour de grand mouvement et doit
remettre son rapport au matin.

\medskip\noindent\textbf{Partie I --- Les chiffres.}
\begin{enumerate}
\item Donner le P\&L réel et l'explication du desk.
\item Donner la part inexpliquée et sa part dans le P\&L.
\item Donner le delta et la vanna du portefeuille.
\item Donner le terme de vanna et le résidu qui en découle.
\item De quoi le résidu est-il constitué~?
\end{enumerate}

\medskip\noindent\textbf{Partie II --- Diagnostic.}
\begin{enumerate}[resume]
\item Quelles sont les trois causes de \omterm{def:rc:pnl-explain-and-independent-price-verification:unexplained}{P\&L inexpliqué} à vérifier en priorité, et dans quel ordre~?
\item Comment l'\omterm{def:rc:pnl-explain-and-independent-price-verification:reval}{explication par réévaluation} localise-t-elle l'effet croisé~?
\item Que montre l'ordre volatilité d'abord~?
\item Pourquoi le \omterm{def:rc:pnl-explain-and-independent-price-verification:unexplained}{P\&L inexpliqué} n'apparaissait-il pas les jours calmes~?
\item Quel rapport la marque en \omterm{def:rc:vanilla-rates-options:flat}{volatilité plate} (flat) du desk entretient-elle avec tout cela~?
\end{enumerate}

\medskip\noindent\textbf{Partie III --- Contrôles.}
\begin{enumerate}[resume]
\item Donner l'ajustement d'IPV et son effet sur la valeur du portefeuille.
\item Donner l'AVA d'incertitude sur le prix de marché, telle que les approches de base et simplifiée la calculeraient.
\item Où placeriez-vous le seuil de \omterm{def:rc:pnl-explain-and-independent-price-verification:unexplained}{P\&L inexpliqué} pour ce desk~?
\item Le desk doit-il être autorisé à expliquer la journée par la seule réévaluation~?
\item Que remonteriez-vous au comité des risques~?
\end{enumerate}

\medskip\noindent\textbf{Partie IV --- Jugement.}
\begin{enumerate}[resume]
\item Quand un \omterm{def:rc:pnl-explain-and-independent-price-verification:unexplained}{P\&L inexpliqué} signale-t-il une erreur de comptabilisation plutôt qu'un risque manquant~?
\item Pourquoi le \omterm{def:rc:pnl-explain-and-independent-price-verification:pc}{product control} doit-il être indépendant du desk~?
\item Comment cela se rattache-t-il au test d'attribution du FRTB du chapitre~23~?
\item Énoncer le \emph{résultat nommé}~: le terme croisé manquant et l'attribution une fois qu'il est ajouté.
\item En une phrase~: que vous indique un \omterm{def:rc:pnl-explain-and-independent-price-verification:unexplained}{P\&L inexpliqué}~?
\end{enumerate}
\end{problem}

\section{Questions d'entretien}

\begin{interviewq}[$\star$ \iqroles{risque, banque}]\label{iq:rc:pnl-explain-and-independent-price-verification:1}
Expliquez le P\&L d'une journée sur un portefeuille d'options. Que faut-il intégrer dans l'explication~?
\end{interviewq}

\begin{interviewq}[$\star\star$ \iqroles{trader, risque}]\label{iq:rc:pnl-explain-and-independent-price-verification:2}
Quelles sont les causes habituelles d'un \omterm{def:rc:pnl-explain-and-independent-price-verification:unexplained}{P\&L inexpliqué}~?
\end{interviewq}

\begin{interviewq}[$\star\star$ \iqroles{banque}]\label{iq:rc:pnl-explain-and-independent-price-verification:3}
Qu'est-ce que la \omterm{def:rc:pnl-explain-and-independent-price-verification:ipv}{vérification indépendante des prix}, et où devient-elle délicate~?
\end{interviewq}

\begin{interviewq}[$\star\star$ \iqroles{chercheur, banque}]\label{iq:rc:pnl-explain-and-independent-price-verification:4}
Qu'est-ce que l'\omterm{def:rc:pnl-explain-and-independent-price-verification:pv}{évaluation prudente}, et en quoi diffère-t-elle de la fair value~?
\end{interviewq}

\begin{interviewq}[$\star\star\star$ \iqroles{développeur}]\label{iq:rc:pnl-explain-and-independent-price-verification:5}
Concevez un système d'explication quotidienne du P\&L pour un grand portefeuille de dérivés.
\end{interviewq}

\begin{interviewq}[$\star\star\star$ \iqroles{risque, banque}]\label{iq:rc:pnl-explain-and-independent-price-verification:6}
Comment détecteriez-vous des valorisations systématiquement favorables à un desk~?
\end{interviewq}
